\section{Discussion and Conclusions}

Evidence for some kind of past inflationary epoch is mounting, although empirical tests have yet to pick out a single, preferred model of inflation from the vast space of candidate models.  So if eternal inflation really were a nearly inevitable implication of any simple inflation mechanism, we would be forced to take very seriously the multiverse conjecture and confront the thorny questions it presents.  Chief among these is the measure problem: how are we to interpret probabilities and make sensible predictions in cosmology if every possible configuration of any subsystem is in principle realized an infinite number of times?

Questions
\begin{itemize}

\item If we assume that the inflaton is a quantum field, will terms in the potential containing odd multiples of the field operator contribute?  Is the state $\Psi$ in which we evaluate $\expv{\phi} = \bra{\Psi}\hat\phi\ket{\Psi}$ one with a definite number of particles?

\item Should we allow the potential to go negative?

\end{itemize}

\subsection{Future Work}

It would be of interest to generalize the numerical methods presented in this work to include multifield models, in which the dynamics are potentially much more complicated.

So why don't we go

somewhere only $W^{\pm}$